\documentclass[12pt,a4paper]{article}
% Compile with XeLaTeX. Requires Amiri, Latin Modern and polyglossia.
\usepackage{fontspec}
\usepackage{amsmath,amssymb}
\usepackage[margin=25mm]{geometry}
\usepackage{polyglossia}
\setdefaultlanguage{arabic}
\setotherlanguage{english}
\setmainfont{Latin Modern Roman}
\newfontfamily\arabicfont[Script=Arabic]{Amiri}
\title{البناء التكراري للأعداد الطبيعية الموجبة من الوحدة المرجعية}
\author{}
\date{\textenglish{26/11/2026} — نسخة \textenglish{0.1}}
\begin{document}
\maketitle
\begin{center}
مسودة تأسيسية للتحكيم الرياضي
\end{center}
\section{نطاق الوثيقة}
الهدف بناء الأعداد الطبيعية الموجبة انطلاقًا من وحدة مرجعية مختارة وعملية تكرار متراكم، دون افتراض العمليات الحسابية المعتادة أو الصفر. هذه وثيقة للنواة المعتمدة، وليست ادعاءً بالأسبقية أو بتأسيس الرياضيات كلها.
\section{المسلمات}
\subsection*{\textenglish{M1} — الوحدة المرجعية}
تُختار مرساة، ويُختار متميز عنها، وتُختار العلاقة بينهما لتكون الوحدة المرجعية، التي يُعبَّر عنها بالعدد \(1\).

نرمز إلى المرساة بـ \(A\)، وإلى المتميز بـ \(B\)، وإلى الوحدة المرجعية بـ \(u\). لا يُفترض أن المرساة عدد، ولا أن العلاقة من نوع هندسي مخصوص.
\subsection*{\textenglish{M2} — التكرار}
يمكن تكرار الموجود مع الاحتفاظ بالأصل وبجميع المخرجات المتراكمة، وتظل إمكانية التكرار قائمة دون حد نهائي.

في الحد الأدنى اللازم للبناء، ينتج التكرار تحققًا جديدًا للوحدة المرجعية، مماثلًا للأصل بالنسبة إلى المرجع المختار ومتميزًا عنه بوصفه تحققًا. لا يحل الجديد محل الموجود، ولا تُدمج التحققات المتمايزة في تحقق واحد. يمكن إعادة تطبيق القاعدة على الوحدة مع استمرار تراكم التحققات السابقة.

هذه التفاصيل بيان للمعنى المقصود بالتكرار، وليست ادعاءً بأنها نتائج مشتقة منطقيًا من عبارة مختصرة.
\section{تعريف المجال المبني}
الأعداد الطبيعية الموجبة هي التعبيرات العددية عن الكميات التي يمكن الحصول عليها بتكرار الوحدة المرجعية تكرارًا متراكمًا منتهيًا، مع اعتبار المخرجات المتكافئة في تعدد التحققات ممثلة للعدد نفسه.

يُستبعد من هذا المجال ما لا يُبنى بهذه الطريقة. هذا حصر تعريفي للمجال، ولا يقتضي وجود مخرج لانهائي.
\section{تمثيل بنائي أولي}
تُعرض التحققات برمز مشترك، ولا يدل تكرار الرمز على تطابق التحققات بوصفها كائنات منفردة:
\[
\begin{array}{c|c}
[u]&1\\
[u,u]&2\\
[u,u,u]&3\\
[u,u,u,u]&4\\
\vdots&\vdots
\end{array}
\]
قاعدة التمثيل: إرفاق تحقق جديد للوحدة بالمخرج الموجود. لا يفترض هذا الإرفاق تعريف الجمع أو الضرب. والأقواس وسيلة تمثيل للتعدد، لا اتحاد مجموعات يمحو التكرار، ولا ادعاء بأن ترتيب التحققات جزء من الدلالة الكمية.

التمثيل يمكن تحقيقه بالكلمات المنتهية غير الفارغة على رمز واحد. هذا نموذج داخل رياضيات قائمة، وليس برهانًا على الاستقلال التأسيسي عن جميع المفاهيم العددية في اللغة الواصفة.
\section{توسع الإمكانات التشغيلية}
لا يُستخدم مفهوم عددي لتحديد مرات التكرار قبل ظهوره في النظام. بعد تحقق كمية جديدة وظهور مفهومها العددي، يصبح استعمالها موضوعًا للتكرار أو استعمال عددها محددًا لمراته ممكنًا، متى توفرت المدخلات. التكرار المتعدد ليس شرطًا لبناء الأعداد؛ يكفي التكرار البسيط المتراكم للوحدة.

يُترك التحديد الرسمي للتكرار المعمم إلى دراسة لاحقة؛ ولا يُعتمد هنا قانون حسابي له.
\section{حدود ما تدعيه الوثيقة}
\begin{itemize}
\item يقدم النموذج مسارًا بنائيًا غير محدود للكميات الموجبة المنتهية.
\item لم تُثبت بعد استقلالية المضمونين المسلّميين أو أدنى عدد ممكن من الافتراضات في لغة صورية محددة. دمجهما في جملة واحدة لا يقلل بالضرورة مضمون الافتراضات.
\item لم يُحسم هل تعريف التعدد والتكافؤ والانتهاء يحتاج إلى بنية أولية إضافية، ولا هل المرجعية العلائقية ضرورية صوريًا للبناء أو مجرد مصدر لاختيار الوحدة.
\item لا يُدّعى أن تسمية الكمية تكفي بذاتها لتعريف العدد تعريفًا صوريًا غير دائري.
\item الصفر والعمليات الحسابية خارج نطاق هذه النسخة؛ لم تُحقق الأسبقية العلمية.
\end{itemize}
\section{طلب التحكيم}
المطلوب نقد رياضي مستقل وصارم: تحديد الغموض أو الدور التعريفي أو الافتراضات المستترة، واختبار الاتساق والكفاية وعدم التكرار غير الضروري للمسلمات، وتقديم نماذج مضادة عند الإمكان. يُفصل بين صلاحية النموذج التمثيلي وصلاحية التأسيس المستقل، وتُحدد أقل التعديلات اللازمة دون تغيير النواة اعتباطًا. لا يُطلب من الناقد الدفاع عن التصور أو إثبات جدته مسبقًا.
\end{document}
